\documentclass[11pt]{article}
\usepackage{rabbithole}

\phasenumber{8}
\phasetitle{Remotes and big repositories}
\phasesubtitle{Refspecs, worktrees, submodules, and what to do when a clone takes an hour.}

\hypersetup{pdftitle={Rabbit Holes, Phase 8: Remotes and big repositories},
            pdfauthor={Accelerate, Manipal University Jaipur},
            pdfsubject={git and GitHub handbook}}

\begin{document}
\maketitlepage

\section{A remote is just a URL with a nickname}

In the workshop you used \co{origin} and \co{upstream} without much explanation.
There is nothing special about either name. A remote is a nickname for a URL, plus
some default rules about which branches map to which.

\begin{shell}
git remote -v
git remote add fork https://github.com/someone/project.git
git remote rename origin github
git remote set-url origin git@github.com:me/project.git
git remote remove old-thing
\end{shell}

You can have as many as you like. Working on someone's pull request often means
adding their fork as a remote, fetching, and checking out their branch, which is
exactly what \co{gh pr checkout} automates for you.

\section{Refspecs}

This is the piece of git that looks like line noise and explains a surprising amount
of behaviour once you can read it.

\begin{shell}
git config --get remote.origin.fetch
\end{shell}

\begin{terminal}
+refs/heads/*:refs/remotes/origin/*
\end{terminal}

Read it as \co{<source>:<destination>}. It says: take everything under
\co{refs/heads/} on the remote, and store it locally under
\co{refs/remotes/origin/}. The leading \co{+} means allow non-fast-forward updates.

This is why \co{origin/main} exists as a thing separate from \co{main}. They are two
different refs. \co{main} is your branch. \co{origin/main} is your local record of
where \co{main} was on the server the last time you fetched. \co{git fetch} updates
the second and never touches the first, which is precisely why fetching is always
safe.

You can write refspecs by hand:

\begin{shell}
git push origin feature:main
git push origin :old-branch
git fetch origin main:local-copy-of-main
\end{shell}

The second one pushes nothing to \co{old-branch}, which deletes it. That is the
original way to delete a remote branch, and \co{git push origin --delete old-branch}
is the readable alias for it.

\section{Tracking branches}

\begin{shell}
git branch -vv
\end{shell}

\begin{terminal}
* main    a1b2c3d [origin/main] Add the cheatsheet
  feature 8e1f4a7 [origin/feature: ahead 2, behind 1] Parser work
\end{terminal}

The bracketed part is the upstream and how far you have diverged. ``Ahead 2, behind
1'' means you have two commits they do not, and they have one you do not. That is
exactly the situation that produces a merge or a rebase.

\begin{shell}
git branch -u origin/main
git push -u origin feature
\end{shell}

\co{-u} sets the upstream, which is what lets you later type bare \co{git push} and
\co{git pull}.

\begin{note}
``Behind'' is measured against your last fetch, not against reality. If you have not
fetched in a week, git will happily tell you that you are up to date. \co{git fetch}
is what makes the answer current, and it changes nothing else, so there is no reason
not to run it often.
\end{note}

\section{Fetch, pull, and the pull that surprises people}

\co{git pull} is \co{git fetch} followed by \co{git merge}. When you have local
commits and the remote has moved, that merge creates a merge commit, and you end up
with ``Merge branch main of github.com...'' scattered through your history.

\begin{shell}
git config --global pull.rebase true
\end{shell}

Now \co{git pull} fetches and rebases instead, replaying your local commits on top of
the remote ones. Your history stays linear and those noise commits disappear.

\begin{trap}
Since Git 2.27, running \co{git pull} with diverged branches and no configured
preference prints a long warning and refuses. The three options it offers are
\co{pull.rebase true}, \co{pull.rebase false}, and \co{pull.ff only}.

Pick \co{true} for personal branches. Pick \co{only} if you would rather git refuse
and let you decide each time, which is the most conservative choice and a reasonable
one.
\end{trap}

Cleaning up references to branches that no longer exist on the server:

\begin{shell}
git fetch --prune
git config --global fetch.prune true
\end{shell}

Without this, \co{git branch -r} accumulates every branch that was ever deleted, and
after a few months of an active project the list is mostly ghosts.

\section{Worktrees}

You are halfway through a change and need to look at another branch. The usual
answer is to stash, switch, look, switch back, and pop.

Worktrees are better. They let one repository have several working directories, each
on a different branch, all sharing one \co{.git}.

\begin{shell}
git worktree add ../project-hotfix hotfix/urgent
git worktree add ../project-review -b review-pr-42
git worktree list
git worktree remove ../project-hotfix
\end{shell}

You now have two directories on disk, on two branches, both live. Your editor can
have both open. Nothing was stashed and nothing was interrupted.

\begin{why}
The disk cost is only the working files, because the history is shared. For a repo
with a long history and a small tree, a second worktree is nearly free.

The pattern this is best at: reviewing someone's pull request while your own work
stays exactly as you left it, including uncommitted changes and a running dev
server.
\end{why}

\begin{trap}
You cannot check out the same branch in two worktrees at once. Git refuses, and it
is right to, because two directories editing one branch would be chaos.
\end{trap}

\section{Big repositories}

When a clone takes too long or fills the disk, there are four separate tools and
they solve different problems.

\subsection{Shallow clone}

Truncate the history.

\begin{shell}
git clone --depth 1 https://github.com/big/project.git
\end{shell}

You get the latest commit and nothing before it. Fast, small, and the standard
choice in CI, where you almost never need history.

Deepen later if you need to:

\begin{shell}
git fetch --deepen 50
git fetch --unshallow
\end{shell}

\begin{trap}
\co{git log} is nearly empty, \co{git blame} is useless, and \co{git bisect} cannot
work. Shallow clones are for machines, not for the copy you develop in.
\end{trap}

\subsection{Partial clone}

Get all the history, but download file contents only when actually needed.

\begin{shell}
git clone --filter=blob:none https://github.com/big/project.git
\end{shell}

You get every commit and tree, so \co{log}, \co{blame} and \co{bisect} all work.
File contents are fetched on demand when you check out or diff.

For most large repositories this is a better default than a shallow clone, and it is
newer, so it is less widely known.

\subsection{Sparse checkout}

Get the whole history but only put part of the tree on disk. Useful in a monorepo
where you work on one service out of forty.

\begin{shell}
git clone --filter=blob:none --sparse https://github.com/big/monorepo.git
git sparse-checkout set services/parser libs/common
git sparse-checkout list
git sparse-checkout disable
\end{shell}

\subsection{Maintenance}

Git can keep a large repository fast in the background.

\begin{shell}
git maintenance start
\end{shell}

This registers scheduled jobs that prefetch, repack, and keep the commit graph
current. On a large repository the difference to \co{git log} and \co{git status}
speed is noticeable.

\section{Submodules}

A submodule is a pointer to a specific commit in another repository, stored inside
yours.

\begin{shell}
git submodule add https://github.com/other/lib.git vendor/lib
git clone --recurse-submodules https://github.com/me/project.git
git submodule update --init --recursive
git submodule update --remote
\end{shell}

The parent repository records the submodule's URL in \co{.gitmodules} and the exact
commit in its tree. That pinning is the point: everyone gets byte-identical
dependencies.

\begin{trap}
Submodules have a poor reputation and it is earned. The recurring problems:

\begin{itemize}
  \item A plain \co{git clone} leaves submodule directories empty, and the error you
    get later is confusing.
  \item \co{git pull} does not update submodules unless you ask.
  \item It is easy to commit a submodule pointer to a commit that only exists on
    your machine, which breaks the build for everyone else and is baffling to
    diagnose.
  \item Branch switching in the parent does not move the submodule.
\end{itemize}

Turn on the setting that removes most of the surprise:

\begin{shell}
git config --global submodule.recurse true
\end{shell}
\end{trap}

\subsection{Subtree, the alternative}

\begin{shell}
git subtree add --prefix=vendor/lib https://github.com/other/lib.git main --squash
git subtree pull --prefix=vendor/lib https://github.com/other/lib.git main --squash
\end{shell}

Subtree copies the other project's files into yours as ordinary files. Anyone
cloning gets everything with no extra commands and no special knowledge.

The trade: your repository grows, and pushing changes back upstream is more awkward.

\begin{note}
Rough guidance. If contributors need to \emph{modify} the dependency and send
changes back, submodule. If they just need it to be there and work, subtree, or
honestly your language's package manager, which is usually the right answer to this
whole question.
\end{note}

\section{Large files}

Git stores a complete copy of every version of every file. That is fine for source
code, which is small and compresses well, and bad for a 50 MB video edited twelve
times.

Git LFS replaces such files with small text pointers and stores the real contents
separately.

\begin{shell}
git lfs install
git lfs track "*.psd"
git lfs track "*.mp4"
git add .gitattributes
git lfs ls-files
\end{shell}

\begin{trap}
\co{git lfs track} only affects files added \emph{after} you set it up. Files
already in history stay there at full size, and the repository stays large. Fixing
that needs a history rewrite, which is Phase 5.

Set up LFS before the first large commit. GitHub also enforces a hard 100 MB limit
per file and will reject a push that exceeds it, which is a rude way to discover
this.
\end{trap}

\section{Reference}

\vspace{0.5em}
\begin{tabularx}{\linewidth}{@{}L{6.6cm} X@{}}
\toprule
\rhtablehead{Goal} & \rhtablehead{Command} \\
\midrule
See where a branch is tracking & \co{git branch -vv} \\
Delete a remote branch & \co{git push origin -{}-delete <branch>} \\
Forget deleted remote branches & \co{git fetch -{}-prune} \\
Rebase instead of merge on pull & \co{git config -{}-global pull.rebase true} \\
Second branch checked out at once & \co{git worktree add ../dir <branch>} \\
Fast clone for CI & \co{git clone -{}-depth 1 <url>} \\
Full history, files on demand & \co{git clone -{}-filter=blob:none <url>} \\
Only part of a monorepo & \co{git sparse-checkout set <paths>} \\
Keep a big repo fast & \co{git maintenance start} \\
Clone including submodules & \co{git clone -{}-recurse-submodules <url>} \\
Track a binary type in LFS & \co{git lfs track "*.mp4"} \\
\bottomrule
\end{tabularx}
\vspace{0.5em}

\checkyourself{
\begin{enumerate}
  \item What is the difference between \co{main} and \co{origin/main}?
  \item Read this refspec aloud: \co{+refs/heads/*:refs/remotes/origin/*}
  \item Why is \co{git fetch} always safe?
  \item When would you use a worktree instead of \co{git stash}?
  \item Shallow clone or partial clone, if you need \co{git blame} to work?
\end{enumerate}}

\vspace{1em}
{\color{rhborder}\rule{\linewidth}{0.8pt}}

\textbf{Next:} Phase 9, \emph{Plumbing}, which opens up \co{.git} and shows you what
a commit really is. After it, most of git's behaviour stops needing to be memorised.

\end{document}
